\documentclass[11pt]{article}
\usepackage{amsmath,amssymb,amsthm}
\usepackage[margin=1in]{geometry}
\newtheorem{theorem}{Theorem}
\newtheorem{proposition}{Proposition}
\theoremstyle{definition}
\newtheorem{remark}{Remark}
\DeclareMathOperator{\Imag}{Im}
\title{The two-row $d=4$ fiber law for $b\equiv 2,3\pmod 4$:\\
a central-binomial closed form, the Gegenbauer engine,\\
and why the obstruction is arithmetic}
\author{Clio}
\date{2026-06-08}
\begin{document}
\maketitle

\noindent\textbf{Setup.} For $\lambda=(2m-b,b)$, $s=1+i$, $A=1+su+u^2$, write
$g_b(m)=[u^b]A^m$, $G_b(m)=[u^b]((1-u)A^m)=g_b(m)-g_{b-1}(m)$, and
$I_b(m)=\Imag G_b(m)\in\mathbb Z[m]$. After dividing out the forced roots $m=0,\dots,R$
($R=\lfloor(b-1)/2\rfloor$), the two-row $d=4$ law $G_{(2m-b,b)}(i)=0\iff(2,2)$ is the statement
\[
(\star)\qquad I_b(m)\neq 0\ \text{ for every integer } m\ge b,\quad b\equiv 2,3\!\!\pmod 4 .
\]
The classes $b\equiv 0,1$ are proved. For $b\equiv 2,3$, $I_b(m)\equiv m\binom{m-1}{R}\pmod 2$:
the bare factor $m$ rules out every single-prime certificate.

\section{The engine}
\begin{proposition}\label{rec}
For all $b\ge 2$,
$\;b\,g_b(m)=s(m-b+1)g_{b-1}(m)+(2m-b+2)g_{b-2}(m)$, with $g_0=1,\ g_1=sm$.
\end{proposition}
\begin{proof}
From $A\,(A^m)'=m(s+2u)A^m$ read off $[u^{b-1}]$.
\end{proof}
\begin{proposition}\label{cf}
$\displaystyle g_b(m)=\binom{m}{b}s^b\,{}_2F_1\!\Big(-\tfrac b2,-\tfrac{b-1}2;\,m-b+1;\,-2i\Big)
= C_b^{(-m)}\!\big(-\tfrac s2\big)$,
the Gegenbauer polynomial of degree $b$ in the \emph{parameter} $-m$ at $x=-s/2$.
\end{proposition}
Because the recurrence coefficient $(2m-b+2)/b$ depends on $m$, the $g_b(m)$ are \emph{not}
orthogonal polynomials in $m$; the classical zero-distribution theory does not apply directly.

\section{The half-integer value}
\begin{theorem}\label{A}
For all $b\ge 1$,
\[
G_b\!\Big(\tfrac{2b-1}{2}\Big)=(-1)^b\,\frac{\binom{2b}{b}}{4^b}\,(1-i)^b ,
\qquad\text{hence}\qquad
I_b\!\Big(\tfrac{2b-1}{2}\Big)=(-1)^{b+1}\,\frac{\binom{2b}{b}}{4^b}\,\Imag(s^b),
\]
of magnitude $2^{\lfloor b/2\rfloor}\binom{2b}{b}/4^b$, vanishing iff $4\mid b$.
\end{theorem}
\begin{proof}
By Proposition~\ref{cf} the lower parameter of $g_b$ at $m=\tfrac{2b-1}2$ is
$m-b+1=\tfrac12$, and that of $g_{b-1}$ is $\tfrac32$. With $z=1-i$ (so $z^2=-2i$,
$1+z=2-i$, $1-z=i$) the quadratic transformations
\[
{}_2F_1(a,a+\tfrac12;\tfrac12;z^2)=\tfrac12\big[(1+z)^{-2a}+(1-z)^{-2a}\big],\quad
{}_2F_1(a,a+\tfrac12;\tfrac32;z^2)=\frac{(1+z)^{1-2a}-(1-z)^{1-2a}}{2z(1-2a)}
\]
(at $a=-\tfrac b2$ and $a=-\tfrac{b-1}2$) give
${}_2F_1(\cdots;\tfrac12;-2i)=\tfrac12[(2-i)^b+i^b]$ and
${}_2F_1(\cdots;\tfrac32;-2i)=\dfrac{(2-i)^b-i^b}{2(1-i)b}$.
Using $\binom{(2b-1)/2}{b}=\binom{2b}{b}/4^b$ and $\binom{(2b-1)/2}{b-1}=2b\binom{2b}{b}/4^b$,
write $C=\binom{2b}{b}/4^b$; with $s/(1-i)=i$,
\[
g_b=C s^{b-1}(1+i)\tfrac12\big[(2-i)^b+i^b\big],\qquad
g_{b-1}=C s^{b-1}(1+i)\tfrac12\big[(2-i)^b-i^b\big].
\]
Subtracting, $(2-i)^b$ cancels: $G_b=C s^{b-1}(1+i)i^b=C(si)^b=(-1)^bC(1-i)^b$, since
$si=-(1-i)$. Finally $(1-i)^b=\overline{s^b}$, so $\Imag G_b=(-1)^{b+1}C\,\Imag(s^b)$ and
$|\Imag(s^b)|=2^{\lfloor b/2\rfloor}$ for $b\not\equiv0$, $0$ for $4\mid b$.
\end{proof}
\begin{remark}
Theorem~\ref{A} unifies both regimes: for $4\mid b$ it gives the rational root $\tfrac{2b-1}2$ of
$Q_b$; for $b\equiv2,3$ it shows the nearest half-integer is never a root, with an explicit
central-binomial value. (Verified symbolically $b\le 15$; magnitudes are the OEIS A001790
numerators, $b\le 23$.)
\end{remark}

\section{Arithmetic, not analytic}
\begin{proposition}\label{arith}
For every $b\equiv 2,3\pmod 4$, $b\ge 6$, the real polynomial $I_b$ has a zero in $(b,\,cb^2)$,
$c\approx 0.33$.
\end{proposition}
\begin{proof}
$I_b(b)$ and $I_b(M)$ (for $M$ past the largest root, where $\operatorname{sgn}I_b=\operatorname{sgn}(\Imag s^b/b!)$)
differ in sign for these $b$; a sign change at integer arguments $\ge b$ forces a real root.
\end{proof}
Hence no uniform analytic lower bound $|I_b(m)|>0$ can hold for real $m\ge b$: any log-concavity /
leading-term dominance argument fails exactly where the in-range real roots live, so it can only
reach $m\gtrsim cb^2$, no better than brute force. The genuine content of $(\star)$ --- that the
irrational in-range roots avoid the integers --- is \emph{arithmetic}. This retires the planned
log-concavity route.

\section{Frontier}
\begin{theorem}
For every $b\equiv 2,3\pmod 4$, $3\le b\le 70$, the only rational roots of $I_b$ are integers
$<b$; thus $Q_b$ has no integer root $\ge b$ and the law holds. Each $Q_b$ is irreducible over
$\mathbb Q$. (Exact rational-root enumeration; supplementary $I_b(m)\neq0$ for all integers
$m\in[b,b^2]$, $b\le150$.)
\end{theorem}

\section*{The gap}
$(\star)$ for $b\equiv2,3$ is open: $Q_b$ has no integer root $\ge b$ (empirically irreducible).
Dead, structurally: single-prime Newton, Eisenstein, $2$-adic truncation, covering sets
(06-08 CODE), and now analytic dominance (Prop.~\ref{arith}). Live, global: uniform
irreducibility of the Gegenbauer-parameter family $C_b^{(-m)}(-(1+i)/2)$; a Legendre
($A^{-1/2}=\sum P_n(-s/2)u^n$) reflection identity; or a family-coefficient growth/Perron bound
on rational roots.
\end{document}
